\documentclass[10pt, a4paper, twocolumn]{article}
\usepackage[utf8]{inputenc}
\usepackage{geometry}
\geometry{a4paper, margin=0.75in, columnsep=0.25in}
\usepackage{hyperref}
\usepackage{booktabs}
\usepackage{verbatim}
\usepackage{titlesec}
\usepackage{mathptmx}
\usepackage{courier}

\title{\textbf{A Hierarchical Service-Class OS Architecture for ARM: Balancing Low-Latency and Strict Isolation}}

\author{
    \textbf{Abdelhalim Yasser ِAbdelhalim}\thanks{Personal Portfolio \& Projects: \url{https://github.com/abdelhalimyasser}} \\
    \textit{Faculty of Computers and Artificial Intelligence, Capital University} \\
    \href{mailto:212400555@fci.capu.edu.eg}{\texttt{212400555@fci.capu.edu.eg}} \\
    \href{https://orcid.org/0009-0009-3629-2168}{ORCID: 0009-0009-3629-2168}
}
\date{\today}

\begin{document}

\maketitle

\begin{abstract}
Modern operating systems face a continuous trade-off between performance and strict isolation. While monolithic architectures offer low latency at the expense of a massive trusted computing base, microkernels provide robust fault containment but suffer from IPC overhead. This paper proposes a novel, conceptual service-class operating system architecture tailored for ARM64 processors, specifically targeting consumer devices where multimedia latency and power efficiency are critical. The architecture introduces a minimal privileged core surrounded by four dynamically isolated service layers (L1-L4). By leveraging ARM's hardware mechanisms, the system facilitates "Same-Layer Protected Fast Calls" for near-native function speed without sacrificing memory safety. Furthermore, we introduce the "Elevator"—a capability-gated, Assembly-level fast path allowing critical user-space services, such as an L2 Earliest Eligible Virtual Deadline First (EEVDF) scheduler, to bypass standard hierarchical constraints. This preprint outlines the architectural blueprint, the C/Rust language strategy, and the future evaluation roadmap.
\end{abstract}

\section{Introduction}
The fundamental problem in modern operating system design is not merely choosing between monolithic and microkernel architectures, but rather acknowledging that different system services possess drastically different trust, criticality, and latency requirements. Subjecting all services to a uniform system-wide performance-isolation trade-off results in inefficiencies, particularly on mobile and portable devices powered by ARM architectures.

For instance, high-bandwidth multimedia pipelines (such as camera and microphone services) require zero-copy data transfers and microsecond latency, whereas third-party drivers demand the strongest possible fault containment. Existing differentiated-isolation systems address parts of this problem but often lack a unified, multi-path communication routing mechanism.

In this paper, we propose a flexible, layer-based architecture that explicitly controls the performance-isolation trade-off at the service level. The design comprises a minimal trusted core handling capability enforcement and hardware interrupts, paired with four distinct service layers. We assign critical resource management (L1) and scheduling (L2, featuring an EEVDF policy) to C-based isolated domains for zero-abstraction performance. Conversely, native system services (L3) and external legacy compatibility (L4) are confined within Rust-based domains to ensure memory safety. Central to our contribution is the introduction of dual communication paths: the standard "Stairs" for normal cross-layer calls, and the "Elevator," a privileged, capability-authenticated Assembly fast-gate for urgent operations.

\section{Related Work}
Existing systems separately address differentiated isolation, mixed-criticality scheduling, fast protected compartments, and capability-based communication.

\begin{itemize}
    \item \textbf{HongMeng Microkernel:} Introduces differentiated isolation classes (IC0, IC1, IC2) to balance performance and strong containment. However, it lacks a complete multi-path communication architecture (such as the proposed capability-gated "Elevator").
    \item \textbf{seL4 \& Mixed-Criticality Systems (MCS):} Provides mathematically verified strong isolation and temporal isolation (CPU budgets). It enforces a uniform strong process boundary and does not relax isolation for high-affinity, same-layer services to achieve near-function-call speeds.
    \item \textbf{ERIM:} Achieves fast, in-process isolation using hardware memory protection keys. It serves as an inspiration for the Same-Layer Protected Fast Call mechanism but is purely an isolation mechanism rather than a full OS architecture.
\end{itemize}

\section{Proposed Architecture}
The proposed architecture is organized as a hierarchical set of isolated service layers built around a minimal privileged kernel. The hierarchy represents different classes of trust, criticality, latency sensitivity, and resource guarantees.

\begin{figure*}[htbp]
\centering
\footnotesize
\begin{verbatim}
                         +----------------------+
                         |       ASM + C        |
                         |     MAIN KERNEL      |
                         | traps / IRQ / MMU    |
                         +----------+-----------+
                                    |
                             ELEVATOR / VIP
                              ASM fast gate
                                    |
                     +--------------+--------------+
                     |                             |
             Normal L1/L2 access            Critical L3/L4
                     |                             |

        === L1: CORE / CRITICAL OS SERVICES (C + ASM) ===
            [ Service A ] <--- local call ---> [ Service B ]
            very low latency, core resource services

        === L2: PERFORMANCE / VIRTUALIZATION (Mainly C) ===
            [ Service C ] <--- local call ---> [ Service D ]
            scheduling policy (EEVDF), virtualization

        === L3: NATIVE SYSTEM SERVICES (C / C++ / Rust) ===
            [ Camera ]    <--- local call ---> [ Mic ]
            multimedia, networking, storage

        === L4: EXTERNAL / LEGACY / THIRD-PARTY (Rust/C) ===
            [ Legacy Driver ]  [ POSIX Comp. ]    Abdelhalim Yasser ِAbdelhalim\thanks{Personal Portfolio & Projects: https://github.com/abdelhalimyasser} ↩
    Faculty of Computers and Artificial Intelligence, Capital University ↩
    212400555@fci.capu.edu.eg ↩
    ORCID: 0009-0009-3629-2168

            strongest containment, least privilege
\end{verbatim}
\caption{The Proposed Layered Service Architecture and Communication Paths}
\end{figure*}

\subsection{Main Kernel}
The Main Kernel contains only mechanisms that require the highest level of privilege: exception and trap entry, context switching, MMU page-table primitives, capability enforcement, and secure call gates.

\subsection{Service Layers (L1 - L4)}
\textbf{L1 (Core Services):} Contains small, highly validated, first-party components such as the physical memory allocator and core power management. \\
\textbf{L2 (Performance \& Virtualization):} Hosts the CPU scheduling policy (e.g., EEVDF) and platform abstraction, separating the mechanism (Kernel) from the policy (L2). \\
\textbf{L3 (Native Services):} Contains isolated services like audio, camera, and network stacks, utilizing Rust for memory safety and zero-copy shared buffers for multimedia throughput. \\
\textbf{L4 (Legacy/External):} The boundary for third-party drivers and POSIX compatibility, providing maximum containment.

\section{Communication and Isolation Model}
The architecture defines three distinct communication routing paths:

\begin{enumerate}
    \item \textbf{Same-Layer Protected Calls:} Services in the same layer communicating frequently use a lightweight token validation mechanism to approach function-call-like performance without removing isolation (e.g., utilizing ARM ASID).
    \item \textbf{Cross-Layer Service Calls ("The Stairs"):} Used for normal operations across layers. The hierarchy does not force unnecessary hops; a service directly reaches its authorized target via the shortest safe path.
    \item \textbf{Elevator / VIP Fast Path:} A protected Assembly-level fast path for critical, latency-sensitive communication. Access requires a valid capability/token authorizing the specific target and operation.
\end{enumerate}

\section{Conclusion and Future Work}
In this paper, we presented a hierarchical, layer-based operating system architecture that explicitly controls the performance-isolation trade-off at the service level. By categorizing services based on criticality, communication affinity, and trust, our design offers a flexible communication routing model comprising same-layer protected calls, normal cross-layer calls, and a privileged Elevator fast path. 

Future work will focus on implementing a proof-of-concept prototype on ARM64. The evaluation will prioritize measuring the latency overhead of the capability-gated Elevator and the context-switch efficiency of same-layer protected calls compared to traditional IPC mechanisms.

\end{document}
